% 13-04(13쪽) — y=f(x) 의 그래프(꺾인 선이 (1,0) 에서 올라가 (2,1) 빈 점에서 끊기고 (2,3) 채운 점부터 가로 반직선).
% 스캔 화질이 낮아 TikZ 로 다시 그림. 원문의 빈 점(○)·채운 점(●)·점선·눈금 1, 3, 1, 2 만 옮긴다.
\begin{tikzpicture}[x=0.52cm, y=0.36cm, line width=0.55pt]
  \draw[->, >=stealth, line width=0.7pt] (-2.1,0) -- (3.6,0) node[below=1pt] {$x$};
  \draw[->, >=stealth, line width=0.7pt] (0,-1.1) -- (0,4.3) node[left=1pt] {$y$};
  \draw[densely dashed, line width=0.4pt] (0,3) -- (2,3);
  \draw[densely dashed, line width=0.4pt] (0,1) -- (2,1);
  \draw[densely dashed, line width=0.4pt] (2,3) -- (2,0);
  \draw[line width=0.6pt] (-1.75,2.75) -- (1,0) -- (2,1);
  \draw[line width=0.6pt] (2,3) -- (3.5,3);
  \draw[fill=white] (2,1) circle (2.2pt);
  \fill (2,3) circle (2.2pt);
  \node[left=1pt, font=\small] at (0,1) {$1$};
  \node[left=1pt, font=\small] at (0,3) {$3$};
  \node[below left=-2pt and -4pt, font=\small] at (0,0) {$\pt{O}$};
  \node[below=1pt, font=\small] at (1,0) {$1$};
  \node[below=1pt, font=\small] at (2,0) {$2$};
  \node[anchor=west, font=\small] at (2.15,3.8) {$y=f(x)$};
\end{tikzpicture}
